%============================================================
% Chapter 17: Nim and XOR Games
% STATUS: PLANNED  (generated by tools/outline.py; replace this file when the chapter is written)
%============================================================
\chapter{Nim and XOR Games}
\label{ch:nim}

\begin{tcolorbox}[enhanced, colback=white, colframe=secondaryblue, boxrule=1pt, arc=2mm]
\textbf{\large\textcolor{primaryblue}{A Nim position is lost exactly when the XOR of its piles is zero: Nim, its variants, and why games add by XOR.}}
\end{tcolorbox}

\chapterfacts{1500}{2600}{4}{8}{30}

\begin{plannedbox}[Planned Contents --- this chapter is not written yet]
\textbf{Part A --- Lessons}
\begin{enumerate}[leftmargin=*, label=\textbf{A\arabic*.}, itemsep=1pt]
    \item Nim: the XOR rule and its proof
    \item Staircase Nim, Moore's Nim$_k$ (sums of bits modulo $k+1$) and misere Nim
    \item Sums of games: Grundy values combine by XOR (Sprague--Grundy theory is the math notes)
    \item Computing Grundy values for small games and spotting their bit patterns
    \item Games on arrays without Nim structure are the array book
\end{enumerate}
\textbf{Part B --- Patterns} (with candidate problems, checked when the chapter is written)
\begin{enumerate}[leftmargin=*, label=\textbf{P\arabic*.}, itemsep=1pt]
    \item \textbf{Nim sums}
    \item \textbf{Staircase Nim}
    \item \textbf{Moore's Nim}
    \item \textbf{Grundy values with bit patterns}
\end{enumerate}
\end{plannedbox}
